\documentclass[10pt,a4paper]{article}
\usepackage[a4paper,margin=0.78in,top=0.76in,bottom=0.72in,headheight=13pt,headsep=15pt,footskip=22pt]{geometry}
\usepackage[T1]{fontenc}
\usepackage{newtxtext,newtxmath}
\usepackage{microtype}
\usepackage{array,tabularx,booktabs}
\usepackage{enumitem}
\usepackage{fancyhdr}
\usepackage{titlesec}
\usepackage{lastpage}
\usepackage[hidelinks]{hyperref}

\hypersetup{pdftitle={Grain Quality Analyzer - Project status and technical handoff},pdfsubject={Project status and technical handoff}}

\pagestyle{fancy}
\fancyhf{}
\fancyhead[L]{\footnotesize Grain Quality Analyzer}
\fancyhead[R]{\footnotesize Project status and technical handoff}
\fancyfoot[C]{\small\thepage\ of \pageref*{LastPage}}
\renewcommand{\headrulewidth}{0.35pt}
\renewcommand{\footrulewidth}{0pt}

\titleformat{\section}{\large\bfseries}{}{0pt}{}
\titleformat{\subsection}{\normalsize\bfseries}{}{0pt}{}
\titlespacing*{\section}{0pt}{10pt}{4pt}
\titlespacing*{\subsection}{0pt}{8pt}{3pt}
\linespread{1.04}
\setlength{\parindent}{0pt}
\setlength{\parskip}{3pt}
\setlist[itemize]{leftmargin=18pt,itemsep=2pt,topsep=2pt,parsep=0pt,label=\textbullet}
\setlist[enumerate]{leftmargin=21pt,itemsep=2pt,topsep=2pt,parsep=0pt,label=\arabic*.}
\renewcommand{\arraystretch}{1.08}
\setlength{\tabcolsep}{6pt}
\newcolumntype{Y}{>{\raggedright\arraybackslash}X}
\newcommand{\hdr}[1]{\textbf{#1}}

\begin{document}

% Title and document details
\begin{center}
{\LARGE\bfseries Grain Quality Analyzer\par}
\vspace{3pt}
{\large Project status and technical handoff\par}
\vspace{6pt}
{\small Prepared 3 November 2026 \quad\textbullet\quad Current repository and documentation reviewed\par}
\end{center}
\vspace{4pt}
\hrule height 0.4pt
\vspace{8pt}

\noindent
\textbf{Current status.} A working web application analyzes uploaded or camera-captured rice-sample images and returns an annotated image, object counts, and per-object measurements. Its active analysis is a deterministic, classical computer-vision pipeline with shape-aware splitting. \textbf{It does not currently run a trained ML model or load model weights.} Predicted counts and whole/broken labels have not been established as production-accuracy or regulatory measurements.

\section{Objective and present scope}
The project aims to help users inspect a rice sample from an image: identify rice grains and visible foreign matter, separate touching grains where possible, show the result clearly, and summarize grain geometry and an image-based whole/broken estimate. The current priority is a usable analysis workflow and an auditable operational record; validated grading is still future work.

\subsection{Original problem statement and coverage}
The original requirement is to detect and quantify, from one smartphone image of 1,000 rice grains: (1) broken; (2) damaged/slightly damaged; (3) discolored; (4) chalky; (5) red; (6) dehusked; (7) immature/shrunken/shriveled; (8) sprouted/weevilled; (9) foreign matter; (10) admixture of lower-class grains; (11) length; (12) breadth; (13) length/breadth ratio; and (14) total count.

\begingroup\small
\setlength{\tabcolsep}{5pt}\renewcommand{\arraystretch}{1.08}
\begin{tabularx}{\linewidth}{@{}>{\centering\arraybackslash}p{0.035\linewidth}p{0.31\linewidth}p{0.095\linewidth}Y@{}}
\toprule
\hdr{\#} & \hdr{Required output} & \hdr{Status} & \hdr{Current coverage} \\
\midrule
1 & Broken grains & Partial & Per-image length heuristic: below 75\% of inferred reference; not validated. \\
2 & Damaged/slightly damaged & Not done & No dedicated damage classification. \\
3 & Discolored & Not done & No discoloration class. \\
4 & Chalky & Not done & No chalkiness classifier. \\
5 & Red grains & Not done & Rice is not classified by color/variety. \\
6 & Dehusked & Not done & No husk-state classifier. \\
7 & Immature/shrunken/shriveled & Not done & No semantic label; a short object may only trigger broken heuristic. \\
8 & Sprouted/weevilled & Not done & No sprout or insect-damage classifier. \\
9 & Foreign matter & Partial & Generic color/brightness blobs; no reliable material/type labels. \\
10 & Lower-class admixture & Not done & No rice class, grade, or admixture classifier. \\
11 & Length & Implemented & Per detected rice instance, in pixels. \\
12 & Breadth & Implemented & Per detected rice instance, in pixels. \\
13 & Length/breadth ratio & Implemented & Computed from pixel measurements. \\
14 & Total count & Partial & Predicted rice instances and foreign objects; not validated for 1,000 grains in one image. \\
\bottomrule
\end{tabularx}
\endgroup
\textit{The current system is therefore a narrower counting/inspection prototype; app, logging, and dashboard work does not implement the missing defect or grade classes.}

\section{Current system and method}
\begingroup\small
\setlength{\tabcolsep}{5pt}
\begin{tabularx}{\linewidth}{>{\leavevmode\bfseries}p{0.14\linewidth}Y}
\toprule
\hdr{Layer} & \hdr{Active implementation} \\
\midrule
Browser & React, TypeScript, Vite, Tailwind; upload/camera capture, result viewer, object selection and inspector. \\
API & FastAPI; \texttt{POST /api/rice-quality/analyze} accepts a multipart image and returns JSON with data-URL images, counts, object records, and summary. \\
Analysis & \texttt{rice\_analysis\_method/rice\_analysis.py}; OpenCV, NumPy, and scikit-image; invoked by the backend's \texttt{label\_image}. \\
Operations & JSON Lines rotating backend logs; optional Streamlit dashboard reads current, rotated, legacy, or exported logs. \\
\bottomrule
\end{tabularx}
\endgroup

The method works on a copy resized to 40\% of source width and height. It converts BGR to HSV, treats pixels outside a green hue/saturation range as foreground, and separates rice from foreign matter using brightness, saturation, hue, and morphological cleanup. Foreign connected components are retained as objects. Rice components are measured for area, solidity, and minimum-area-rectangle length/width.

An image-specific shape model is learned from isolated, elongated, high-solidity components. It uses the longer isolated grains to estimate a representative whole-grain shape; if too few suitable components exist, an area-based fallback estimates dimensions. Rice components that look like one grain are kept. Larger or irregular clumps are split by fitting mixtures of elongated anisotropic Gaussian shapes, selecting the smallest candidate count that passes a coverage/overshoot criterion, then refining cuts with marker-controlled watershed at image seams. A merge step attempts to reunite over-split pieces and discard crumbs. The random generator has a fixed default seed.

For each rice object, the backend computes length and breadth in source-image pixels. The reference length estimated by the method is scaled back to source coordinates. A grain is called whole at length $\geq 0.75$ of this per-image reference and broken below it. With no valid reference, status is undetermined; it must not be interpreted as whole or broken. This is a heuristic, scale-dependent image measure, not a mass-based or certified grade.

\newpage
\section{Work delivered to date}
\begin{itemize}
  \item \textbf{End-to-end workflow:} Responsive browser upload and camera capture connect to the API; results include source and masked images, counts, IDs, measurement records, and a CSV/image download path.
  \item \textbf{Inspection UI:} Whole grains are blue, broken grains yellow, and foreign objects red; undetermined rice is gray. Users can select an object to inspect its measurements and isolated mask. Image zoom is pointer-centered, pan is optional, and IDs/boxes start hidden.
  \item \textbf{Input and memory controls:} Accepted formats are JPEG, PNG, WebP, TIFF, and BMP; upload cap 25 MiB, request cap 26 MiB, decoded cap 16 megapixels, and minimum side 8 px. The service allows one active analysis and two queued requests; excess load receives HTTP 503. It normalizes EXIF orientation, grayscale, alpha, and 16-bit input for the 8-bit CV path. Working masks are retained as object crops and mapped to source coordinates for response rendering.
  \item \textbf{Privacy-conscious event logging:} Request, analysis, and lifecycle events include IDs, status, durations, dimensions, result counts, queue wait, RSS/CPU samples, and upload byte count. Image bytes and original filenames are excluded. Logs rotate at 10 MiB with five backups; filesystem persistence must be configured on ephemeral hosts.
  \item \textbf{Operations dashboard:} Streamlit presents request/analysis latency, HTTP outcomes, endpoint volume, sampled memory/CPU, upload and queue distributions, prediction-count trends, and log integrity. Graphs are full width with explanatory notes. It explicitly cannot infer model accuracy without reviewed labels.
  \item \textbf{Project cleanup and guidance:} Current application code and the active method are distinguished from archived models, experiments, and the former pipeline. Root and dashboard READMEs describe setup, API, constraints, logging, deployment, graphs, and interpretation.
\end{itemize}

\subsection{Data and evidence available}
The local repository has \texttt{data/rice\_count/raw/} (860 images; local and Git-ignored) and \texttt{data/rice\_count/test/} (20 images; included in the public GitHub repository). The test filenames encode nominal whole, broken, and stone counts, but are not reviewed instance masks or independently verified ground truth. EXIF/comment metadata was removed from the public test copies without changing decoded pixels. Project history says images may show the same grains at the same distance, so exact capture-session and sample provenance should be confirmed before designing a split. No reviewed per-instance masks are currently wired into the dashboard, and the active application does not train on these images.

\textbf{Historical proxy count run (24 October).} The fresh classical proxy processed 860 raw images against user-supplied nominal references of 100 rice grains and 10 foreign objects per image: 92.70\% combined count agreement (6,903 absolute count errors / 94,600 nominal reference objects), rice MAE 7.59, foreign-matter MAE 0.43, and both counts exact in 135/860 images (15.7\%). This is agreement with fixed nominal counts, \textbf{not independent image-specific ground-truth validation}. The run preceded project consolidation; rerun the current active version and confirm labels before using these results for method comparison.

Previous local checks recorded in the work log include running the 20 test images through API checks for output dimensions/bytes, mask bounds, and count/object consistency; comparing image decoding with the prior OpenCV path; and checking TIFF, invalid/oversized/chunked uploads, alpha, EXIF orientation, 16-bit input, CORS, and four-request concurrency. The log records peak sampled RSS of about 153 MiB on that local set. These are historical checks, not a current rerun or proof of detection accuracy or hosted capacity. The committed copies are available under \texttt{data/rice\_count/test/} in the GitHub repository link below; the 860 raw images remain private/local and are not included.

\subsection{Issues encountered and current state}
\begingroup\small
\setlength{\tabcolsep}{5pt}\renewcommand{\arraystretch}{1.08}
\begin{tabularx}{\linewidth}{>{\raggedright\arraybackslash\leavevmode\bfseries}p{0.17\linewidth}YY}
\toprule
\hdr{Issue} & \hdr{What happened} & \hdr{Current interpretation} \\
Old pipeline import & Startup once failed on undefined \texttt{Tuple\_Validation} in a retired grain-profile path. & The active API now imports the current CV method; the old path is archived, not the live analysis route. \\
Missing experimental model & Startup warning reported missing \texttt{backend.models.grain\_infer}. & That experimental module is not required by the active route; this is not evidence of a loaded model. \\
Hosted memory & Render reported the instance exceeded its 512 MiB limit and restarted. & Input caps, serialized analysis/short queue, crop masks, and decoder/memory work were added. Local sampling is encouraging; hosted peak/load behavior remains unverified. \\
\bottomrule
\end{tabularx}
\endgroup

\newpage
\section{Open risks and questions for review}
\begin{enumerate}
  \item \textbf{Ground truth and split design.} Audit how the 860 raw and 20 test images relate by physical grain batch, scene, capture session, and near-duplicate frames. If the same grains appear across images, split by sample/session (not randomly by image) to prevent leakage. Create a small, double-reviewed benchmark with instance masks and rice/foreign labels before comparing methods.
  \item \textbf{Detection and segmentation performance.} Measure instance precision/recall/F1, rice and foreign-object count MAE/bias, and mask IoU or AP against reviewed annotations. Report results by isolated versus touching/overlapping grains, broken fragments, foreign-object type, lighting, and background. Dashboard prediction trends are not accuracy metrics.
  \item \textbf{Broken-grain validity.} The 75\% rule uses a reference inferred from the same image and lengths in pixels. Test its stability across mixture composition, view angle, camera distance, and broken-grain prevalence. If the intended decision is a standards-based or mass-based grade, define the relevant standard, physical calibration, sampling protocol, and mass comparison before presenting this estimate as compliance.
  \item \textbf{Method choice.} Establish a classical-CV baseline on the reviewed benchmark first. Then assess whether a labeled instance-segmentation model is justified by accuracy, label cost, runtime, deployment memory, and maintainability. Treat pseudo-labels from the current method as unverified until human review.
  \item \textbf{Robustness and deployment.} Probe changes in green-sheet color, illumination, shadows, exposure, camera, resolution, overlap, and foreign materials. Load-test the deployed service at the target instance size; measure peak RSS and tail latency under concurrent uploads. Check persistent log retention and avoid storing images unless a reviewed privacy/retention design is approved.
  \item \textbf{Original target scale.} Include a representative single image containing 1,000 grains in the benchmark; current tests and nominal 100/10 count run do not demonstrate performance at this density.
\end{enumerate}

\section{Suggested first work with the expert}
\textbf{1. Agree on the target.} Decide whether success means grain counting, foreign-matter detection, separation of touching grains, broken-grain proportion, or a staged combination. Define acceptable errors and the operating conditions. Define annotation rules for all 14 requested outcomes, especially ambiguous defect classes and what qualifies as lower-class admixture.

\begin{enumerate}[start=2,label=\bfseries\arabic*.,leftmargin=19pt]
  \item Inventory the image set and document sample/session provenance, camera/setup variation, duplicates, and permitted use.
  \item Annotate and adjudicate a representative, leakage-safe benchmark; reserve a genuinely independent test split.
  \item Freeze the current pipeline as a baseline. Measure class-wise detection, count, mask, broken/whole, runtime, and memory outcomes with confidence intervals and error examples.
  \item Use error analysis to choose the next intervention: acquisition controls, color/geometry calibration, revised CV, or a trained instance-segmentation model. Repeat the same evaluation after each change.
  \item Set deployment acceptance thresholds and monitor operational health separately from image-quality performance.
\end{enumerate}

\section{Useful handoff locations}
\begingroup\small
\setlength{\tabcolsep}{5pt}\renewcommand{\arraystretch}{1.12}
\begin{tabularx}{\linewidth}{>{\raggedright\arraybackslash\leavevmode\ttfamily}p{0.40\linewidth}Y}
\toprule
\hdr{Path} & \hdr{Purpose} \\
rice\_analysis\_method/\allowbreak rice\_analysis.py & Active segmentation, shape model, splitting, and output instances. \\
backend/app/api/routes.py & Image decode, limits, CV invocation, object metrics, and API response. \\
frontend/src/ & Upload/camera workflow and results/inspection UI. \\
backend/app/main.py; backend/app/services/\allowbreak event\_logger.py & File logging, request lifecycle and operational events. \\
monitoring\_dashboard/app.py; monitoring\_dashboard/README.md & Log parsing, metrics, graphs, and interpretation caveats. \\
data/rice\_count/\{raw,test\}/ & Local image sets; verify provenance and annotations before evaluation. \\
progress/WORK\_LOG.md & Dated change history and recorded checks/limitations. \\
\bottomrule
\end{tabularx}
\endgroup

\vfill
\hrule height 0.4pt
\vspace{3pt}
{\footnotesize\textit{This handoff describes the repository's active implementation and recorded work as reviewed on 3 November 2026. Validate predictive performance on a leakage-safe, independently annotated test set before operational or scientific use.}}

\section*{Project links}
\begin{description}[leftmargin=1.55in,labelwidth=1.45in,itemsep=2pt,topsep=3pt]
  \item[Backend API (Render):] \url{https://grain-quality-analyzer.onrender.com}
  \item[Frontend (Vercel):] \url{https://grain-quality-analyzer.vercel.app/}
  \item[Source code (GitHub):] \url{https://github.com/kiran-8287/grain-quality-analyzer}
  \item[Demo video:] \textit{[Insert demo video URL here]}
\end{description}
\end{document}
